\documentclass[11pt]{article}
\usepackage[margin=2.4cm]{geometry}
\usepackage[T1]{fontenc}
\usepackage[brazilian,shorthands=off]{babel}
\usepackage{amsmath,amssymb}
\usepackage{booktabs}

\title{Estudo: redes neurais recorrentes}
\author{Italo Brito Brandão}
\date{18/10/2026}

\begin{document}
\maketitle

\section{A ideia}

Uma rede feedforward trata cada entrada como um evento isolado. Uma RNN (rede neural recorrente) lê uma sequência e carrega um estado \(h_t\) de um passo para o próximo. Esse estado é a memória do que já foi visto.

O uso típico é texto, áudio e série temporal: a ordem dos elementos importa, e o comprimento da sequência não é fixo.

\section{A recorrência}

No passo \(t\) a rede recebe o vetor \(x_t\) e o estado anterior \(h_{t-1}\):

\begin{equation}
  h_t = \tanh(W_{xh} x_t + W_{hh} h_{t-1} + b_h)
  \label{eq:rnn}
\end{equation}

\begin{equation}
  y_t = W_{hy} h_t + b_y
\end{equation}

\(W_{xh}\) mistura a entrada nova. \(W_{hh}\) mistura a memória. A mesma matriz \(W_{hh}\) é usada em todos os passos: a rede não ganha um peso novo para cada posição da sequência.

\(h_0\) costuma ser o vetor zero. A saída \(y_t\) pode ser usada em todo passo (etiquetar cada token) ou só no último (classificar a sequência inteira).

\section{Exemplo numérico}

Uma RNN minúscula, com estado e entrada escalares. Pesos \(W_{xh} = 1\), \(W_{hh} = 0{,}5\), vieses zero, e a sequência \(x = (1, 0, 1)\).

\begin{center}
\begin{tabular}{cccc}
  \toprule
  \(t\) & \(x_t\) & conta interna & \(h_t = \tanh(\cdot)\) \\
  \midrule
  1 & 1 & \(1\cdot 1 + 0{,}5\cdot 0 = 1\) & \(\tanh 1 \approx 0{,}76\) \\
  2 & 0 & \(0 + 0{,}5\cdot 0{,}76 = 0{,}38\) & \(\tanh 0{,}38 \approx 0{,}36\) \\
  3 & 1 & \(1 + 0{,}5\cdot 0{,}36 = 1{,}18\) & \(\tanh 1{,}18 \approx 0{,}83\) \\
  \bottomrule
\end{tabular}
\end{center}

No passo 2 a entrada é zero, mas \(h_2\) não é zero: o \(1\) do primeiro passo ainda está na memória, atenuado por \(0{,}5\). No passo 3 os dois sinais se somam.

\section{Treino}

O treino desenrola a recorrência no tempo e aplica backpropagation através desses passos (BPTT). O gradiente de \(W_{hh}\) é um produto de matrizes, um fator por passo. Se os valores singulares de \(W_{hh}\) ficam menores que 1, o produto vai a zero: a rede esquece o começo da sequência. Se ficam maiores que 1, o gradiente explode.

Por isso uma RNN simples, a da equação~\eqref{eq:rnn}, aprende dependências curtas e falha em dependências longas.

\section{LSTM e GRU}

LSTM e GRU trocam o \(\tanh\) puro por portões que decidem o que guardar e o que descartar. O estado deixa de ser multiplicado por \(W_{hh}\) em todo passo, e o gradiente consegue atravessar muitos timesteps.

\begin{center}
\begin{tabular}{ll}
  \toprule
  modelo & o que muda \\
  \midrule
  RNN & um estado, atualizado sempre \\
  GRU & portões de atualização e de esquecimento, um estado \\
  LSTM & portões de entrada, esquecimento e saída, mais uma célula \(c_t\) \\
  \bottomrule
\end{tabular}
\end{center}

Na prática, texto longo hoje usa Transformer. RNN, LSTM e GRU continuam úteis quando a sequência chega um elemento por vez e o estado precisa ser pequeno: áudio em tempo real, séries, e modelos embarcados.

\section{O que observar ao implementar}

\begin{itemize}
  \item O batch é uma matriz de sequências. Passos de tempo andam no eixo da sequência, e \(h\) anda junto.
  \item Sequências do mesmo batch têm comprimentos diferentes. O padding não pode entrar na perda nem no estado.
  \item Explosão de gradiente aparece como loss virando \texttt{nan}. O primeiro controle é cortar a norma do gradiente.
  \item Desaparecimento de gradiente aparece como a rede ignorando o início da frase. Aí o modelo simples já não basta: GRU ou LSTM.
\end{itemize}

\end{document}
